%% ma: L11-B01-0006
%% lop: 11
%% chuong: 1
%% bai: 1
%% dang: 11.1.7
%% loai: TL
%% muc: TH
%% dap_an: "-"
%% nguon: "de-cuong-giua-ky-1-toan-11-nam-2026-2027-truong-thpt-hoang-van-thu-ha-noi-2-8.pdf (M08, Đề cương giữa kì 1 Toán 11 THPT Hoàng Văn Thụ, Hà Nội), Phần 4 Bài 2 ý 2"
%% kien_thuc: "Bài 1 (giá trị lượng giác của góc lượng giác)"
%% trang_thai: chinh-thuc
%% ghi_chu: "Đã duyệt 29/10/2026."
\begin{ex}
Cho tam giác $ABC$, chứng minh rằng $\sin(A+B+2C)=-\sin C$.
\loigiai{Vì $A+B+C=\pi$ nên $A+B+2C=\pi+C$. Do đó $\sin(A+B+2C)=\sin(\pi+C)=-\sin C$.}
\end{ex}
